%!TEX root main.tex

% general macros
\newcommand{\tuple}[1]{(#1)}
\newcommand{\tuplesmall}[1]{(#1)}
\newcommand{\goesto}[1]{\xrightarrow{#1}}
\newcommand{\from}{\leftarrow}
\newcommand{\e}{\varepsilon}
\newcommand{\sep}{\mid}

\newcommand{\ignore}[1]{}

% English
\newcommand{\wlg}{w.l.o.g.}
\newcommand{\Wlg}{W.l.o.g.}
\newcommand{\wrt}{w.r.t.}
\newcommand{\vs}{vs.}
\newcommand{\cf}{cf.}
\newcommand{\eg}{e.g.}
\newcommand{\Eg}{E.g.}
\newcommand{\aka}{a.k.a.}
\newcommand{\ie}{i.e.}
\newcommand{\Ie}{I.e.}
\newcommand{\rhs}{r.h.s.}
\newcommand{\lhs}{l.h.s.}
\newcommand{\egs}{EGS}
\newcommand{\etal}{\emph{et al.}}
\newif\ifstartedinmathmode
\newcommand*{\st}{%s.t.~
  \relax\ifmmode\startedinmathmodetrue\else\startedinmathmodefalse\fi
  \ifstartedinmathmode{\;\cdot\;}\else{s.t.}\fi%
}

\makeatletter
\providecommand*{\cupdot}{%
  \mathbin{%
    \mathpalette\@cupdot{}%
  }%
}
\newcommand*{\@cupdot}[2]{%
  \ooalign{%
    $\m@th#1\cup$\cr
    \hidewidth$\m@th#1\cdot$\hidewidth
  }%
}
\makeatother

% sets
\newcommand{\A}{\mathbb A}
\renewcommand{\AA}{\mathcal A}
\newcommand{\B}{\mathbb B}
\newcommand{\BB}{\mathcal B}
\newcommand{\C}{\mathbb C}
\newcommand{\D}{\mathbb D}
\newcommand{\F}{\mathbb F}
\renewcommand{\H}{\mathbb H}
\newcommand{\N}{\mathbb N}
\newcommand{\Z}{\mathbb Z}
\newcommand{\Q}{\mathbb Q}
\newcommand{\R}{\mathbb R}
\renewcommand{\S}{\mathbb S}
\newcommand{\Sleft}{\S^{\mathsf L}}
\newcommand{\Sright}{\S^{\mathsf R}}
\newcommand{\Ssyn}{\S_{\text{syn}}}
\renewcommand{\SS}{\mathcal S}
\newcommand{\T}{\mathbb T}
\newcommand{\Tleft}{\T^{\mathsf L}}
\newcommand{\Tright}{\T^{\mathsf R}}
\newcommand{\Tsyn}{\T_{\text{syn}}}
\newcommand{\TT}{\mathcal T}
\newcommand{\V}{\mathbb V}
\renewcommand{\k}{\kappa}
\renewcommand{\a}{\mathfrak a}

\newcommand{\x}{\mathtt x}
\newcommand{\y}{\mathtt y}
\newcommand{\z}{\mathtt z}

\newcommand{\FF}{\mathcal F}
\newcommand{\GG}{\mathcal G}

% maths
\let\oldcirc\circ
\let\circ\relax
\DeclareMathOperator{\circ}{\oldcirc\,}

\newcommand{\compose}[1]{\,\oldcirc_{#1}\,}

\DeclareMathOperator{\binconv}{\ast}
% \DeclareMathOperator{\rel}{{\ensuremath R}}
\DeclareMathOperator{\bisimulation}{\sim}

\newcommand{\rel}{\mathrel R}
\newcommand{\relation}[1]{\mathrel {#1}}
\newcommand{\tilderel}{\mathrel {\widetilde R}}
\newcommand{\tilderelation}[1]{\mathrel {\widetilde {#1}}}

% METAFONT mumbo-jumbo for shuffle
% \DeclareSymbolFont{Shuffle}{U}{shuffle}{m}{n}
% \DeclareFontFamily{U}{shuffle}{}
% \DeclareFontShape{U}{shuffle}{m}{n}{%
%   <-8>shuffle7%
%   <8->shuffle10%
% }{}
% \DeclareMathSymbol\shuffle{\mathbin}{Shuffle}{"001}
% \DeclareMathSymbol\cshuffle{\mathbin}{Shuffle}{"002}

% \let\oldshuffle\shuffle
% \let\shuffle\relax
% \DeclareMathOperator{\shuffle}{\oldshuffle}
% \newcommand{\shufflepower}[2]{{#1}^{\oldshuffle#2}}
% \newcommand{\shinv}[1]{\shufflepower {#1}{-1}}

% shuffle by hand
% https://tex.stackexchange.com/questions/183988/how-to-make-the-shuffle-product-symbol-in-latex-math-mode
\makeatletter
\providecommand*{\shuffle}{%
  \mathbin{\mathpalette\shuffle@{}}%
}
\newcommand*{\shuffle@}[2]{%
  % #1: math style
  % #2: unused
  \sbox0{$#1\vcenter{}$}%
  \kern .15\ht0 % side bearing
  \rlap{\vrule height .25\ht0 depth 0pt width 2.5\ht0}%
  \raise.1\ht0\hbox to 2.5\ht0{%
    \vrule height 1.75\ht0 depth -.1\ht0 width .17\ht0 %
    \hfill
    \vrule height 1.75\ht0 depth -.1\ht0 width .17\ht0 %
    \hfill
    \vrule height 1.75\ht0 depth -.1\ht0 width .17\ht0 %
  }%
  \kern .15\ht0 % side bearing
}
\makeatother

\let\oldcolon\colon
\let\colon\relax
\DeclareMathOperator{\colon}{\,:\,}

\let\oldmod\mod
\let\mod\relax
\DeclareMathOperator{\mod}{mod}

\DeclareMathOperator{\cauchy}{\ast}
\DeclareMathOperator{\hadamard}{\odot}
\DeclareMathOperator{\infiltration}{\uparrow}
\DeclareMathOperator{\product}{\ast}
% \DeclareMathOperator{\sq}{\Box}%{\circ_\C}
% \DeclareMathOperator{\binconv}{\ast}
\newcommand{\norm}[2]{\left| #2 \right|_{#1}}
\newcommand{\onenorm}[1]{\norm 1 {#1}}
\newcommand{\inftynorm}[1]{\norm \infty {#1}}
\newcommand{\height}[1]{\inftynorm{#1}}
\newcommand{\limplies}{\Rightarrow}
\newcommand{\valname}{v}
\newcommand{\val}[1]{\valname({#1})}
\newcommand{\valpname}[1]{\valname_{#1}}
\newcommand{\valp}[2]{\valpname {#1}({#2})}
\newcommand{\valstar}[1]{\valname^*({#1})}
\newcommand{\eval}[2]{\left.#1\right|_{#2}}
\newcommand{\floor}[1]{\lfloor#1\rfloor}
\newcommand{\abs}[1]{\left|#1\right|}
\newcommand{\card}[1]{\left|#1\right|}
\newcommand{\length}[2]{\left|#2\right|_{#1}}
\newcommand{\len}[1]{\left|#1\right|}
\newcommand{\set}[1]{\left\{#1\right\}}
\newcommand{\setof}[2]{\set{#1 \;\middle|\; #2}}
\newcommand{\multiset}[1]{\left\{\!\left\{#1\right\}\!\right\}}
\newcommand{\multisetof}[2]{\multiset{#1 \mid #2}}
\newcommand{\census}[1]{\mathsf{c}\left(#1\right)}
\newcommand{\eqclass}[2]{\left[#2\right]_{#1}}
% \newcommand{\averagefun}{\mathsf{avg}}
% \newcommand{\average}[2]{\averagefun_{#1}\left(#2\right)}
\newcommand{\average}[2]{\left\{#2\right\}_{#1}}
\newcommand{\commclass}[1]{\widetilde {#1}}
\newcommand{\rad}[1]{\sqrt{#1}}
\newcommand{\polyexpr}[2]{#1\{#2\}}
\newcommand{\poly}[2]{#1[#2]}
\newcommand{\polyof}[3]{\poly {#1} {#2 \mid #3}}
\newcommand{\diffpoly}[2]{\poly{#1}{#2}^{(*)}}
\newcommand{\ncpoly}[2]{#1\langle#2\rangle}
\newcommand{\rational}[2]{#1(#2)}
\newcommand{\Lie}[2]{\mathsf{Lie}_{#1}\langle#2\rangle}
\newcommand{\LieDeg}[3]{\mathsf{Lie}^{#1}_{#2}\langle#3\rangle}
\newcommand{\rat}[2]{#1(#2)}
\newcommand{\lparenangle}{{\,\mathclap{(}\!\langle}}
\newcommand{\rparenangle}{{\,\mathclap{\rangle}\!)}}
\newcommand{\ncrat}[2]{#1\lparenangle#2\rparenangle}
\newcommand{\nclaurent}[2]{#1\lparenangle\!\lparenangle#2\rparenangle\!\rparenangle}
\newcommand{\powerseries}[2]{#1[\![#2]\!]}
\newcommand{\algps}[2]{\powerseries {#1} {#2}_{\mathsf{alg}}}
\newcommand{\laurentseries}[2]{#1(\!(#2)\!)}
\newcommand{\series}[2]{#1\langle\!\langle#2\rangle\!\rangle}
\newcommand{\lcseries}[2]{#1\langle\!\langle#2\rangle\!\rangle^{\text{LC}}}
\newcommand{\commseries}[2]{#1\langle\!\langle#2\rangle\!\rangle^{\mathsf{comm}}}
\newcommand{\algseries}[2]{#1\langle\!\langle#2\rangle\!\rangle_{\mathsf{alg}}}
\newcommand{\sequences}[2]{#1[\![\N^{#2}]\!]}
\newcommand{\matrices}[3]{#3^{#1 \times #2}}
% \newcommand{\lin}[2]{\langle #2 \rangle_{#1}}
% \newcommand{\linof}[3]{\lin {#1} {#2 \mid #3}}
\newcommand{\ideal}[1]{\langle#1\rangle}
\newcommand{\idealof}[2]{\ideal{#1 \mid #2}}
\newcommand{\diffideal}[1]{\ideal{#1}^{(*)}}
\newcommand{\diffidealof}[2]{\idealof{#1}{#2}^{(*)}}
\newcommand{\module}[1]{\langle#1\rangle}
\newcommand{\moduleof}[2]{\module{#1 \mid #2}}
\newcommand{\lin}[2]{\mathsf{Lin}_{#2}\left(#1\right)}
% \newcommand{\linof}[3]{\lin{#1 \mid #2} {#3}}
\renewcommand{\dim}[1]{\mathsf{dim}\,{#1}}
\newcommand{\leadingterm}[1]{\mathsf{lt}\,{#1}}
\newcommand{\sem}[1]{\llbracket#1\rrbracket}
\newcommand{\Asem}[2]{#1\sem{#2}}
\newcommand{\support}[1]{\mathsf{supp}(#1)}
% \newcommand{\lang}[1]{L\left(#1\right)}
\newcommand{\bigO}[1]{O\left(#1\right)}
\newcommand{\bigOmega}[1]{\Omega\left(#1\right)}
\newcommand{\multiplicity}[2]{\#({#2})_{#1}}
\newcommand{\Parikh}[1]{\mathrm c({#1})}
\newcommand{\Parikhinv}[1]{\mathrm c^{-1}({#1})}
\newcommand{\coefficient}[2]{[#1]#2}
\newcommand{\restrict}[2]{\left.#1\right|_{#2}}
\newcommand{\restrictsmall}[2]{#1|_{#2}}
\newcommand{\equivmod}[3]{{#1} \equiv {#2}\; (\mathsf{mod}\ #3)}
\newcommand{\ord}[1]{\mathsf{ord}(#1)}
\newcommand{\comm}[2]{\mathsf{comm}_{#1}(#2)}
\newcommand{\stops}[1]{\mathsf{s2p}\left(#1\right)}
\newcommand{\pstos}[1]{\mathsf{p2s}\left(#1\right)}
\newcommand{\iso}{\cong}
\newcommand{\completion}[1]{\widehat{#1}}
\newcommand{\linearspan}[2]{\mathsf{Lin}_{#1}(#2)} %{\langle#1\rangle}
\newcommand{\inner}[2]{\langle#1, #2\rangle}
\newcommand{\closure}[1]{\overline{#1}}
\newcommand{\reversal}[1]{{#1}^\mathsf{rev}}
\newcommand{\transpose}[1]{{#1}^\mathsf{T}}
\newcommand{\id}{\mathsf{id}}
%\newcommand{\leftright}[3]{{#2}^{#1, #3}}
\newcommand{\leftright}[3]{\mbox{}^{#1}{#2}^{#3}}
\newcommand{\lr}[3]{\leftright{#1}{#2}{#3}}
\renewcommand{\iff}{\Leftrightarrow}

\newcommand{\seq}[2]{\left(#1 \mapsto {#2}\right)}
\newcommand{\derive}[1]{\delta_{#1}}
\newcommand{\deriveleft}[1]{\delta^L_{#1}}
\newcommand{\deriveright}[1]{\delta^R_{#1}}
\DeclareRobustCommand{\atled}{\text{\reflectbox{$\delta$}}}
\newcommand{\derleft}[2]{{#1}^{-1} {#2}}
\newcommand{\derright}[2]{{#2} {#1}^{-1}}
\newcommand{\class}[2]{[#1]_{#2}}
\let\oldpartial\partial
% \renewcommand{\partial}[2]{\frac {\oldpartial {#2}} {\oldpartial {#1}}}
\renewcommand{\partial}[2]{\oldpartial_{#1}{#2}}
\newcommand{\partialk}[3]{\oldpartial_{#1}^{#2}{#3}}
\newcommand{\bigpartial}[2]{\frac {\oldpartial {#2}} {\oldpartial{#1}}}
\newcommand{\dn}[2]{\mathrm d^{#1} {#2}}
\renewcommand{\d}[1]{\dn {} {#1}}
\newcommand{\total}[2]{\frac {\d {#1}}{\d {#2}}}
\newcommand{\totaln}[3]{\frac {\dn {#1} {#2}}{\d {#3}^{#1}}}
\newcommand{\jacobian}[2]{\partial{#1}{#2}}
\newcommand{\shift}[1]{\sigma_{#1}}
\newcommand{\deductionrule}[2]{\begin{array}{c}%
{#1} \\ \hline {#2}%
\end{array}}
\newcommand{\one}{\mathbbm 1}
\newcommand{\zero}{\mathbb 0}

\newcommand{\X}{\mathcal X}
\newcommand{\Y}{\mathcal Y}
\newcommand{\ZZ}{\mathcal Z}
\newcommand{\ZERO}{\mathbf 0}
\newcommand{\ONE}{\mathbf 1}
\newcommand{\EGS}[1]{\mathsf{EGS}[#1]}
\newcommand{\SETSP}{\mathsf{SET}}
\newcommand{\SET}[1]{\SETSP[#1]}
\newcommand{\SEQSP}{\mathsf{SEQ}}
\newcommand{\SEQ}[1]{\SEQSP[#1]}
\newcommand{\CYCSP}{\mathsf{CYC}}
\newcommand{\CYC}[1]{\CYCSP[#1]}
\newcommand{\TREE}[2]{\mathsf{T}_{#2}[#1]}
\newcommand{\CATALAN}[1]{\mathsf{T}[#1]}
\newcommand{\CAYLEY}[1]{\mathsf{C}[#1]}
\newcommand{\SPGRAPHS}[1]{\mathsf{SP}[#1]}
\newcommand{\SGRAPHS}[1]{\mathsf{S}[#1]}
\newcommand{\PGRAPHS}[1]{\mathsf{P}[#1]}

\newcommand{\lneg}{\neg}
\newcommand{\erf}{\mathrm{erf}}
\newcommand{\arsinh}{\mathrm{arsinh}}
\newcommand{\arcosh}{\mathrm{arcosh}}
\newcommand{\artanh}{\mathrm{artanh}}
\newcommand{\csch}{\mathrm{csch}}
\newcommand{\sech}{\mathrm{sech}}

% acronyms
\newcommand{\CDF}{\textsf{CDF}}
\newcommand{\CDA}{\textsf{CDA}}
\newcommand{\BPP}{\textsf{BPP}}
\newcommand{\WBPP}{\textsf{WBPP}}
\newcommand{\WFA}{\textsf{WFA}}
\newcommand{\WTS}{\textsf{WTS}}
\newcommand{\LTS}{\textsf{LTS}}
\newcommand{\ATS}{\textsf{ATS}}

\newcommand{\NC}{\textsf{NC}}
\newcommand{\coRNC}{\textsf{coRNC}}
\newcommand{\PTIME}{\textsf{PTIME}}
\newcommand{\coRP}{\textsf{coRP}}
\newcommand{\PSPACE}{\textsf{PSPACE}}
\newcommand{\EXPTIME}{\textsf{EXPTIME}}
\newcommand{\NEXPTIME}{\textsf{NEXPTIME}}
\newcommand{\coNEXPTIME}{\textsf{coNEXPTIME}}
\newcommand{\EXPSPACE}{\textsf{EXPSPACE}}
\newcommand{\TWOEXPTIME}{{2-\EXPTIME}}
\newcommand{\TWONEXPTIME}{{2-\NEXPTIME}}
\newcommand{\coTWONEXPTIME}{{2-\coNEXPTIME}}
\newcommand{\TWOEXPSPACE}{{2-\EXPSPACE}}

\newcolumntype{A}{
 >{$}r<{$}
 @{\extracolsep{0pt}}
 >{${}} l <{$}
% @{\extracolsep{\fill}}
} % A for "align"
%% (1) "r" column in math mode:          >{$} r <{$}
%% (2) no space:                         @{}
%% (3) "l" column in math mode, with
%%     an empty subformula at the start: >{${}} l <{$}

\newcolumntype{M}{
 >{$}c<{$}
}

% theorems
\theoremstyle{plain}
\newtheorem{theorem}{Theorem}
\newtheorem*{theorem*}{Theorem}
\newtheorem{lemma}[theorem]{Lemma}
\newtheorem{corollary}[theorem]{Corollary}
\newtheorem{claim}[theorem]{Claim}
\newtheorem*{claim*}{Claim}
\newtheorem{proposition}[theorem]{Proposition}
\newtheorem*{proposition*}{Proposition}
\newtheorem{fact}[theorem]{Fact}
\newtheorem*{fact*}{Fact}

\theoremstyle{remark}
\newtheorem{remark}{Remark}
\newtheorem{example}{Example}

\theoremstyle{definition}
\newtheorem{definition}{Definition}

\theoremstyle{problem}
\newtheorem{problem}{Problem}

\let\oldproof\proof
\let\endoldproof\endproof

% \makeatletter
% \let\proof\@undefined
% \let\endproof\@undefined 
% \makeatother
% \newenvironment{proof}{\oldproof}{\endoldproof}

% \newenvironment{myproof}{\comment}{\endcomment}
\newenvironment{myproof}{\proof}{\endproof}

\newenvironment{claimproof}{\begin{proof}[Proof of the claim]}{\end{proof}}

% \newenvironment{myclaimproof}{\comment}{\endcomment}
\newenvironment{myclaimproof}{\begin{claimproof}}{\end{claimproof}}

\crefname{equation}{}{}%{Eq.}{Eqns.}
\crefname{definition}{Definition}{Definitions}
\crefname{claim}{Claim}{Claims}
\crefname{proposition}{Proposition}{Propositions}
\crefname{lemma}{Lemma}{Lemmas}
\crefname{corollary}{Corollary}{Corollaries}
\crefname{theorem}{Theorem}{Theorems}
\crefname{figure}{Fig.}{Figures}
\crefname{fact}{Fact}{Facts}
\crefname{example}{Example}{Examples}
\crefname{remark}{Remark}{Remarks}
\crefname{section}{§}{§§}
\crefname{subsection}{§}{§§}
\crefname{subsubsection}{§}{§§}
\crefname{enumi}{}{}
\crefname{problem}{Problem}{Problems}

\crefname{table}{Table}{Tables}
\crefname{figure}{Figure}{Figures}

\crefname{appendix}{Appendix}{Appendices}

% TODO
\newcommand{\ilorenzo}[1]{\todo[inline, color=blue!30]{\textsf L: #1}}
\newcommand{\lorenzo}[1]{\todo[color=blue!30]{\textsf L: #1}}